\section{Deferred Constructions}\label{app:deferred-constructions}

This appendix collects the technical details deferred from the main text. Each
subsection below records the statement to be proved, the expected witness
objects, the proof shape, and the soundness/completeness points that should be
filled in.

\subsection{Batched Table Lookup}\label{app:concatenated-table-lookup}

Given content $\mathbf{F}_i\in \mathbb{F}^{n \times m_i}$
  and table $\mathbf{T}_i\in \mathbb{F}^{N\times m_i}$ for $i\in [L]$,
  and commitments to their corresponding polynomial encodings
  $F_i(X)$ over $\mathbb{H}_{n m_i}$
  and $T_i(X)$ over $\mathbb{H}_{N m_i}$,
  our goal is to prove
  \[
    \mathbf{F}_0\| \cdots \mathbf{F}_{L-1} \subseteq \mathbf{T}_0\| \cdots \mathbf{T}_{L-1},
  \]
  where notation ``$\subseteq$'' means
  that each row of the left matrix appears as a row of the right matrix.

The key challenge here is that the table widths $m_i$ vary
  by nature.
  For example, consider a one-row matrix $(a, b)$.
  Suppose it is committed as $f(X)$ over $\mathbb{H}_{1\times 2}$,
  its corresponding vector over $\mathbb{H}_{1\times 4}$ is
  \[
    \mathbf{f}' = (a, \square, b, \square).
  \]
  Note that $f(\omega_4^0) = f(\omega_2^0) = a$
  and $f(\omega_4^2) = f(\omega_2^1) = b$.
  However, there is no guarantee what values
  $f(\omega_4^1)$ and $f(\omega_4^3)$ will be.

Our insight is hence to create a mask
  \[
      \mathbf{m} = (1, 0, 1, 0),
  \]
  so $f(X)m(X)$ honestly commits to $(a, 0, b, 0)$ over $\mathbb{H}_{1\times 4}$.
  By \Cref{lem:row-replication}, $m(X) = \lambda^{(0)}_2(X^2)$.

Formally, let $m = \max_i m_i$. We pad each matrix $\mathbf{F}_i$
  to $\widehat{\mathbf{F}}$ of width $m$ by
  \begin{align}
    & \widehat{F}(X) = F(X)\cdot \lambda^{(0)}_{m/m_i}(X^{nm_i}),\\
    & \text{where } \lambda^{(0)}_{m/m_i}(X) = \frac{X^{m/m_i} - 1}{m/m_i(X-1)}.
  \end{align}
  And we pad each table $\mathbf{T}_i$ to $\widehat{\mathbf{T}}_i$ of width $m$ similarly.

Now that the contents and the tables are aligned,
  we can combine them into a single table lookup by a verifier challenge $\alpha$:
  \begin{align}
    \sum_i \alpha^i\cdot \widehat{\mathbf{F}}_i \subseteq \sum_i \alpha^i\cdot \widehat{\mathbf{T}}_i.
  \end{align}
  This is reduced to a table lookup argument in \cite{cq++24}.

We note that this construction also supports preprocessing,
  at the cost of computing cached quotient~\cite{cq22}
  for all $\widehat{\mathbf{T}}_i$ in $O(Lm\log m)$.

\subsection{Column-Replication}\label{app:replication}
Given a vector $\mathbf{f}\in\mathbb{F}^n$ and
  its committed polynomial encoding $f(X)$ and a positive integer $m$,
  we want to prove
  that committed polynomial $G(X)$ encodes the matrix
  \[
    \mathbf{G} = (\mathbf{f}\,\|\,\mathbf{f}\,\|\,\cdots\,\|\,\mathbf{f}) \in \mathbb{F}^{n\times m}.
  \]
  Unlike \Cref{lem:row-replication}, we do not know if there
  is a closed-form expression for $G(X)$ in terms of $f(X)$.
  So below we give a direct proof of correctness of $G(X)$.

We prove the following two identities:
\begin{align}
    &\forall i\in[n], j\in[m-1].\ \mathbf{G}[i, j+1] - \mathbf{G}[i, j] = 0,\\
    &\forall i\in [n].\ \mathbf{G}[i, 0] = \mathbf{f}[i],
\end{align}
which are proven separately by
\begin{align}
  \exists q_1(X).&\
  G(X) - G(\omega_{nm}^{-1}X) = q_1(X)\cdot \nu_n(X),\\
  \exists q_2(X).&\
  G(X) - f(X) = q_2(X)\cdot \nu_n(X),
\end{align}
i.e. two zero-checks.

\subsection{Concrete Constructions for Segmented Prefix Sum}\label{app:prefix-sum}

We now give the concrete constructions deferred from
\Cref{sec:zkgbdt-protocol}. Throughout, we write
$[t]:=\{0,1,\dots,t-1\}$.
For a tensor $\mathbf{F}\in\mathbb{F}^{a\times b\times c}$,
we need the following three relations.

\paragraph{Target constraints.}
Given a commitment to $\mathbf{F}$, prove commitments to $\mathbf{G}$ and
$\mathbf{H}$ such that:
\begin{enumerate}
    \item \textbf{Prefix along the last dimension.}
    \[
        \mathbf{G}[i,j,k]
        =
        \sum_{k'=0}^{k}\mathbf{F}[i,j,k'],
        \quad
        \mathbf{H}[i,j]
        =
        \sum_{k'\in[c]}\mathbf{F}[i,j,k'].
    \]

    \item \textbf{Prefix along the middle dimension.}
    \[
        \mathbf{G}[i,j,k]
        =
        \sum_{j'=0}^{j}\mathbf{F}[i,j',k],
        \quad
        \mathbf{H}[i,k]
        =
        \sum_{j'\in[b]}\mathbf{F}[i,j',k].
    \]

    \item \textbf{Prefix along the first dimension.}
    \[
        \mathbf{G}[i,j,k]
        =
        \sum_{i'=0}^{i}\mathbf{F}[i',j,k],
        \quad
        \mathbf{H}[j,k]
        =
        \sum_{i'\in[a]}\mathbf{F}[i',j,k].
    \]
\end{enumerate}

We note that the first and third types of prefix sums
  are actually prefix sums along the row and column directions of a matrix, respectively.
  For notation simplicity, we write
  matrix prefix sums for these two types in the following.
\paragraph{Matrix row sum / row-wise prefix sum.}
Let $\mathbf{F}\in\mathbb{F}^{n\times m}$ be the committed matrix, and let
$\mathbf{H}\in\mathbb{F}^{n}$ be the claimed row sum, i.e.
\begin{align}
    \mathbf{H}[i] = \sum_{j\in[m]} \mathbf{F}[i,j].
\end{align}
To prove this relation, we first extend $\mathbf{H}$ to a matrix
$\widehat{\mathbf{H}}\in\mathbb{F}^{n\times m}$ by padding zeros before the
last column:
\begin{align}
    \widehat{\mathbf{H}} = (\mathbf{0}\,\|\,\mathbf{H}) \in \mathbb{F}^{n\times m}.
\end{align}
Let
\[
    \mathbb{H}_{n}^{\mathrm{last}}
    :=
    \omega_{nm}^{m-1}\mathbb{H}_n
\]
be the last-column subdomain, and let
$\nu_{\mathbb{H}_{n}^{\mathrm{last}}}(X)$ denote its vanishing polynomial.
The relation between $\mathbf{H}$ and $\widehat{\mathbf{H}}$ is enforced by
proving
\begin{align}
    &\exists q_1(X).\ 
    \widehat{H}(X)
    =
    q_1(X)\cdot
    \frac{\nu_{nm}(X)}{\nu_{\mathbb{H}_{n}^{\mathrm{last}}}(X)},\\
    &\exists q_2(X).\
    \widehat{H}(X)-H(\omega_{nm}^{-(m-1)}X)
    =
    q_2(X)\cdot \nu_{\mathbb{H}_{n}^{\mathrm{last}}}(X),
\end{align}
which are reduced to two zero-checks.
Intuitively, the first condition says that $\widehat{\mathbf{H}}$ is zero
outside the designated positions, while the second says that on those positions
it agrees with $\mathbf{H}$ after shifting the last-column domain back to
$\mathbb{H}_n$.

Next, we commit to an auxiliary witness
$\mathbf{B}\in\mathbb{F}^{n\times m}$, defined as the row-wise prefix sum of
$\mathbf{F}-\widehat{\mathbf{H}}$:
\[
    \mathbf{B}[i,j]
    :=
    \sum_{k=0}^{j}\bigl(\mathbf{F}[i,k]-\widehat{\mathbf{H}}[i,k]\bigr).
\]
We then prove the following two identities:
\begin{align}
  &\exists q_3(X).\
  B(X) = q_3(X)\cdot \nu_{\mathbb{H}_{n}^{\mathrm{last}}}(X),\\
  &B(X)-B(\omega_{nm}^{-1}X)=F(X)-\widehat{H}(X).
\end{align}
The first identity guarantees that $\mathbf{H}$ is indeed the correct row sum,
while the second guarantees that $\mathbf{B}$ is the correct prefix sum witness.

The key observation is that $\mathbf{B}$ is the de facto prefix sum of
$\mathbf{F}-\widehat{\mathbf{H}}$:
\begin{align}
    \mathbf{B}[i,j]
    =
    \underbrace{\mathbf{B}[i-1\bmod n,m-1]}_{=0}
    +
    \sum_{k=0}^{j}\bigl(\mathbf{F}[i,k]-\widehat{\mathbf{H}}[i,k]\bigr).
\end{align}
Hence $\mathbf{B}[i,m-1]=0$ implies
\begin{align}
  \sum_{k\in[m]}\mathbf{F}[i,k]=\mathbf{H}[i].
\end{align}


To obtain the full row-wise prefix relation, note that
\begin{align}
  \mathbf{G} = \mathbf{B} + \widehat{\mathbf{H}}.
\end{align}

\paragraph{The second type prefix sum.}
We now describe how to prove the second type relation
\begin{align}
    \mathbf{G}[i,j,k]
    =
    \sum_{j'=0}^{j}\mathbf{F}[i,j',k],
    \qquad
    \mathbf{H}[i,k]
    =
    \sum_{j'\in[m]}\mathbf{F}[i,j',k],
\end{align}
where $\mathbf{F}\in\mathbb{F}^{n\times m\times \ell}$,
$\mathbf{G}\in\mathbb{F}^{n\times m\times \ell}$, and
$\mathbf{H}\in\mathbb{F}^{n\times \ell}$.
The difficulty is that, under the row-major flattening of
$\mathbf{F}$, the entries contributing to one fixed $\mathbf{H}[i,k]$ are not
stored contiguously: they appear once every $\ell$ positions inside the
$i$-th block of length $m\ell$.

\smallskip
\noindent
\textbf{Step 1: construct a mask.}
We first commit to a mask matrix
\begin{align}
    \mathbf{M}\in\mathbb{F}^{n\times m\times \ell},
\end{align}
such that
\begin{align}
    \mathbf{M}_{\{n\times m\times \ell\}}[i,j,k]
    =
    \begin{cases}
        1, & j=m-1,\\
        0, & \text{otherwise}.
    \end{cases}
\end{align}
Thus, within each block of $m$ rows corresponding to a fixed $i$, the mask is
nonzero exactly on the last row. Equivalently, the support of $\mathbf{M}$
marks the positions at which the full sum over the middle dimension should
appear.

To prove that $\mathbf{M}$ is correct, we prove
\begin{align}
    M(X)-M(\omega_{nm\ell}^{-1}X)
    =
    \lambda_{m\ell}^{(0)}(X^n\omega_{nm}^{m-1})-\lambda_{m\ell}^{(0)}(X^n).
\end{align}
The right-hand side is explicit and can be
computed by the verifier. This identity uniquely characterizes the desired
periodic mask.

\smallskip
\noindent
\textbf{Step 2: construct an index witness.}
Next, we commit to an index matrix
\begin{align}
    \mathbf{I}\in\mathbb{F}^{n\times m\times \ell},
\end{align}
defined as the (vector) prefix sum of the mask multiplied by the mask:
\begin{align}
  &\mathbf{I} = \mathbf{I}'\odot \mathbf{M},\\
  &\text{where } \mathbf{I}'_{\{nm\ell\}}[i] = \sum_{i'=0}^{i}\mathbf{M}_{\{nm\ell\}}[i'].
\end{align}
$\mathbf{I}$ will be used as an indicator of the positions 
where the sum should be copied to.

\smallskip
\noindent
\textbf{Step 3: lift $\mathbf{H}$ to the shape of $\mathbf{F}$.}
We then commit to an aligned tensor
\[
    \widehat{\mathbf{H}}\in\mathbb{F}^{n\times m\times \ell},
\]
such that
\[
    \widehat{\mathbf{H}}[i,j,k]
    =
    \begin{cases}
        \mathbf{H}[i,k], & j=m-1,\\
        0, & \text{otherwise}.
    \end{cases}
\]
In other words, $\widehat{\mathbf{H}}$ stores each $\mathbf{H}[i,k]$ exactly at the
positions selected by the mask, and is zero elsewhere.

The relation between $\mathbf{H}$ and $\widehat{\mathbf{H}}$ is proved through a copy
constraint guided by the index witness $\mathbf{I}$. Concretely, after
flattening $\mathbf{H}$ as a vector of length $n\ell$, we prove
\begin{align}
    (1,2,\dots,n\ell)^\top \,\|\, \mathbf{H}
    \subseteq
    \mathbf{I} \,\|\, \widehat{\mathbf{H}},\label{eq:copy-constraint-H-Hhat}
\end{align}
where $\subseteq$ denotes a lookup. Intuitively, this states that whenever
$\mathbf{I}$ indicates the $t$-th target location, the value written in
$\widehat{\mathbf{H}}$ at that location is exactly the $t$-th entry of $\mathbf{H}$.

\smallskip
\noindent
\textbf{Step 4: prove the sum by a column-prefix argument.}
Finally, we consider
\[
    \mathbf{R}:=\mathbf{F}-\widehat{\mathbf{H}}.
\]
We prove that the column-wise prefix sum of $\mathbf{R}$ vanishes at all marked
positions. More precisely, let
\begin{align}
    \mathbf{P}_{\{nm\times \ell\}}[i, j]
    :=
    \sum_{j'=0}^{j}\mathbf{R}_{\{nm\times \ell\}}[i, j'].\label{eq:column-prefix-vanishes-at-mask}
\end{align}
The correctness of $\mathbf{P}$ follows from the proof of the first type of prefix sum.

We show that $\mathbf{P}$ vanishes at all positions where $\mathbf{M}$ is nonzero, i.e.
\begin{align}
  \mathbf{P}\odot \mathbf{M} = \mathbf{0}.
\end{align}

This finishes the proof of the correctness of $\mathbf{H}$.
  To prove the correctness of $\mathbf{G}$, we can simply prove $\mathbf{G} = \mathbf{P} + \widehat{\mathbf{H}}$.

\smallskip
\noindent
\textbf{Soundness and completeness.}
Soundness follows from the soundness of the primitives.
To see the completeness, note that
by applying copy constraint (\Cref{eq:copy-constraint-H-Hhat}),
we enforce that the only nonzero term in $\widehat{\mathbf{H}}$ is
$\mathbf{H}[i,k]$ at $j=m-1$, so \Cref{eq:column-prefix-vanishes-at-mask} is equivalent to
\[
    \sum_{j'=0}^{m-1}\mathbf{F}[i,j',k]-\mathbf{H}[i,k]=0,
\]
namely,
\[
    \mathbf{H}[i,k]=\sum_{j\in[m]}\mathbf{F}[i,j,k].
\]

\begin{example}\label{ex:second-prefix-sum}
Consider $\mathbf{F}\in\mathbb{F}^{2\times 3\times 4}$. After flattening it as
a matrix of size $6\times 4$, the entries of
$\mathbf{H}\in\mathbb{F}^{2\times 4}$ should appear at the last row of each
block of three rows. Thus the mask has the form
\[
\mathbf{M}=
\begin{pmatrix}
0&0&0&0\\
0&0&0&0\\
1&1&1&1\\
0&0&0&0\\
0&0&0&0\\
1&1&1&1
\end{pmatrix}.
\]
Over $\mathbb{H}_{nm\ell}$, $M(X) - M(\omega_{nm\ell}^{-1}X)$ encodes
\[
\begin{pmatrix}
-1&0&0&0\\
0&0&0&0\\
1&0&0&0\\
-1&0&0&0\\
0&0&0&0\\
1&0&0&0
\end{pmatrix} = 
\begin{pmatrix}
0&0&0&0\\
0&0&0&0\\
1&0&0&0\\
0&0&0&0\\
0&0&0&0\\
1&0&0&0
\end{pmatrix} - 
\begin{pmatrix}
1&0&0&0\\
0&0&0&0\\
0&0&0&0\\
1&0&0&0\\
0&0&0&0\\
0&0&0&0
\end{pmatrix}
\]
The $\mathbf{I}$ and $\mathbf{I}'$ used in the proof are respectively
\[
\mathbf{I}=
\begin{pmatrix}
0&0&0&0\\
0&0&0&0\\
1&2&3&4\\
0&0&0&0\\
0&0&0&0\\
5&6&7&8
\end{pmatrix},
\quad
\mathbf{I}'=
\begin{pmatrix}
0&0&0&0\\
0&0&0&0\\
1&2&3&4\\
4&4&4&4\\
4&4&4&4\\
5&6&7&8
\end{pmatrix}.
\]
\end{example}
